\documentclass[11pt,compress,t,notes=noshow, aspectratio=169, xcolor=table]{beamer}

\usepackage{../../style/lmu-lecture}
% Defines macros and environments
\input{../../style/common.tex}

\title{Applied Machine Learning}
% \author{LMU}
%\institute{\href{https://compstat-lmu.github.io/lecture_iml/}{compstat-lmu.github.io/lecture\_iml}}
\date{}

\begin{document}
\lecturechapter{Ensembling: 1 Level Stacking}
\lecture{Applied Machine Learning}

% ------------------------------------------------------------------------------

\begin{vbframe}{Stacking}
Train multiple base models on the same data and use their predictions as input features for a meta model that makes the final ensemble prediction.
\end{vbframe}

% ------------------------------------------------------------------------------

\begin{vbframe}{Steps}

\begin{itemize}
\item Base model training
\item Prediction Collection
\item Meta model training
\end{itemize}
\end{vbframe}

% ------------------------------------------------------------------------------

\begin{vbframe}{Base Model Training}
\end{vbframe}

% ------------------------------------------------------------------------------

\begin{vbframe}{Collect Predictions}
Build task for the meta model
Note on using original features or not
Cross-validated predictions
Overfitting
\end{vbframe}

% ------------------------------------------------------------------------------

\begin{vbframe}{Meta Model Training}
\end{vbframe}

% ------------------------------------------------------------------------------

\begin{vbframe}{Popular Meta Models}
  \begin{itemize}
    \item linear / logistic regression
    \item decision tree
    \item random forest
    \item ...
  \end{itemize}
\end{vbframe}

% ------------------------------------------------------------------------------


\begin{vbframe}{}
  Note on interpretability when using interpretable meta models

\begin{itemize}
  \item ...
\end{itemize}

\end{vbframe}


% ------------------------------------------------------------------------------

\endlecture
\end{document}
